\section{Relationships, cooperation, and conflict}\label{sec:paul-relations}

Paul wrote with an emphatic first-person voice, but seldom as a man without
partners. Letters name co-senders, emissaries, hosts, patrons, household
churches, laborers, prisoners, competitors, and people he had wounded or who
had wounded him. His apostolic identity was tested in these relationships.
Communion never meant an absence of conflict; conflict never proves that every
relationship ended in schism.

\subsection{Cephas, James, John, and Jerusalem}

After the interval he calls three years, Paul went to Jerusalem to visit
Cephas and stayed fifteen days; he also saw James, the Lord's brother. The
visit joined a newly called apostle to the principal Easter witnesses without
making Paul their delegate. At a later meeting James, Cephas, and John,
recognized as pillars, perceived the grace given to Paul and Barnabas and
extended the right hand of fellowship. Their differentiated missions toward
the circumcised and Gentiles did not divide two gospels; Paul immediately says
they shared concern for the poor.

Antioch exposed a serious failure in that fellowship. Paul says Cephas had
eaten with Gentiles but withdrew when people came from James, fearing the
circumcision group. Other Jewish believers, including Barnabas, joined the
withdrawal. Paul opposed Cephas publicly because table separation contradicted
the truth of the gospel. Galatians records no reply, reconciliation scene, or
permanent rupture. First Corinthians later treats Cephas as a recognized
apostle and rejects a Corinthian faction using either apostle's name. Roman
memory eventually pairs their witness.

The dispute concerned the terms of Gentile incorporation in Christ and what
shared meals embodied. Paul did not require Titus or Gentile believers
generally to be circumcised. Acts says he did circumcise Timothy, whose mother
was Jewish, in a different pastoral setting. Whether every detail can be
synchronized is less important than preserving the different questions.

Paul's relationship with Jerusalem remained costly. He defended his mission's
independence, gathered aid, carried the collection, accepted a purification
proposal in Acts, and repeatedly sought unity. He also faced suspicion from
believers and lethal opposition. Neither romantic harmony nor an invention of
two unrelated religions does justice to the sources.

\subsection{Barnabas, Silas, Timothy, and Titus}

Barnabas appears in Acts as Paul's advocate, teacher, co-delegate, and first
major mission partner; Galatians presents him as known to the churches and
trusted in Jerusalem. His withdrawal at Antioch hurt precisely because the
partnership mattered. The later disagreement over Mark was sharp enough to
separate their routes. No source gives Paul the last word on Barnabas's
subsequent fruitfulness.

Silvanus or Silas shares the address of both Thessalonian letters and travels
with Paul in Acts. Timothy is co-sender of six undisputed or disputed letters,
Paul's envoy to anxious communities, and one whose welfare mattered intensely.
Titus, absent from Acts by name, is indispensable in Galatians and Second
Corinthians: an uncircumcised Greek whose status tests Gentile freedom, a
delegate in a damaged Corinthian relationship, and an organizer of the
collection. The difference between Acts and the letters is a warning against
equating narrative prominence with historical importance.

Mark is the object of the Acts dispute with Barnabas. Later canonical letters
associated with Paul speak positively of Mark, but Colossians and Second
Timothy have disputed authorship. Philemon, an undisputed letter, includes
Mark among Paul's co-workers and securely prevents the disagreement from being
turned into proof of lifelong enmity.

\subsection{Women as patrons, apostles, and laborers}

Paul's undisputed letters name women in the mission with concrete verbs and
offices. Phoebe is a deacon of the church at Cenchreae and a benefactor or
patron of many, including Paul; Romans commends her to its recipients and
therefore makes her a plausible bearer and interpreter of the letter, though
the delivery role is an inference. Prisca or Priscilla is named before Aquila
in several greetings, hosts a house church, risks her life, and is called a
co-worker. Junia is a fellow Jew and prisoner, described with Andronicus as
prominent among the apostles or well known to them; the grammar and scope are
debated, but the ancient female name should not be silently masculinized.

Mary, Tryphaena, Tryphosa, Persis, and Julia appear in Romans 16. Euodia and
Syntyche struggled beside Paul in the gospel and are urged to reconciliation.
Chloe's people bring the report that triggers much of First Corinthians.
Apphia shares the address of Philemon. The women who pray and prophesy in
First Corinthians 11 are not hypothetical, even though Paul regulates their
appearance. These witnesses make an all-male picture of Pauline mission
historically untenable.

They also sharpen the hard texts. First Corinthians 14:34--35 tells women to
be silent, creating tension with chapter 11; scholars dispute whether the
verses address a particular speech, were displaced, or are a later
interpolation. First Timothy 2 restricts women through an argument from
creation, but the Pastorals' direct authorship is disputed. Historical honesty
requires both the women whose labor Paul celebrates and the restrictive
reception that Pauline texts have supported.

\subsection{Households, enslaved persons, and patrons}

Assemblies met in homes, making household resources and status unavoidable.
Gaius hosts Paul and the whole church; Stephanas's household devotes itself to
service; Prisca and Aquila host assemblies in more than one city; Philemon and
Apphia belong to a household church. Patronage could protect a mission and
also create expectations Paul resisted. His refusal of Corinthian support and
acceptance of Philippian gifts show context-sensitive judgment, not disdain
for every benefactor.

Enslaved and free persons shared these assemblies inside a slave society.
Philemon concerns Onesimus, whom Paul calls his child and beloved brother. Paul
sends him and asks Philemon to receive him as he would Paul, offering to absorb
any debt. The letter destabilizes possession by invoking brotherhood, but it
does not demand abolition of slavery as an institution, record Onesimus's
manumission, or permit modern readers to ignore the coercive structure.
Later household codes in Colossians and Ephesians address masters and slaves
within that structure; their authorship and reception require separate study.

\subsection{Rivals and wounded communities}

Paul argues against teachers pressing circumcision on Gentiles in Galatia,
against factional leaders at Corinth, and against ``super-apostles'' whom he
believes preach another Jesus. We usually hear the opponents through Paul's
polemic, so the concrete demands and actions he describes are more recoverable
than stable identities or a complete account of their positions.

Paul's language can be fierce. Galatians uses bodily sarcasm; Second
Corinthians attacks rivals' motives; Philippians calls some workers dogs.
Ancient controversy explains rhetoric but does not make every tactic an
imitable pastoral norm. The same letters show Paul apologizing for causing
pain, fearing he has labored in vain, begging reconciliation, and refusing to
lord authority over faith. His holiness lies neither in perpetual
agreeableness nor in sanctifying every angry sentence, but in a ministry
continually summoned back to the cross he proclaimed.
